% Lösungen Einheit 3 von 3 – Mit Parameter (zusammenbau v0.1)
\einheitenkopf{Einheit 3 von 3 · Mit Parameter}

\erg{12}{a) nur im Nenner \quad b) $P = \frac{0{,}6a}{0{,}6a + 0{,}1 \cdot (1 - a)}$; Probe mit $a = 0{,}4$: $P = 0{,}8$ \quad c) Im Term steht b nur im Nenner, beim Pfad der reklamierenden Stammkunden; der Zähler, der Pfad der reklamierenden Neukunden, bleibt gleich. Sinkt b, sinkt der Nenner, also wächst der Bruch \quad d) Graph A: Kauft kein anderer Besucher Eis, kam jeder Eiskäufer mit dem Rad, der Graph beginnt bei eins; kaufen alle anderen Eis, bleibt ein Wert zwischen null und eins – B endet bei null, C beginnt unter eins \quad e) $a \approx 0{,}5$; a ist der Anteil der WLAN-Nutzer unter den privat Reisenden; $f(a) = 0{,}5$ heißt: die Hälfte der WLAN-Nutzer reist geschäftlich (nicht: die Hälfte der Geschäftsreisenden nutzt das WLAN) \quad f) $x \ge 0{,}82$ (aus $\frac{0{,}1(1 - x)}{0{,}1(1 - x) + 0{,}9} \le 0{,}02 \Leftrightarrow 0{,}098(1 - x) \le 0{,}018$; nach oben gerundet)}
\erg{13}{a) Fehler: ein Zahlenbeispiel beweist die Aussage nicht allgemein; richtig: der Zähler ist konstant, der Nenner fällt mit b, also wächst der Bruch für jedes fallende b \quad b) Man teilt dieselbe Größe durch eine kleinere Zahl; jeder Anteil wird größer, also wächst der Quotient \quad c) $w \le 0{,}0096$ (aus $\frac{0{,}0095}{0{,}0095 + 0{,}99w} \ge 0{,}5 \Leftrightarrow 0{,}99w \le 0{,}0095$); der Test darf also höchstens etwa jeden hundertsten Gesunden falsch positiv melden}
